% Project:  TeX PhD Motivation Letter Template
% Author:   Deep Ghuge

% EDIT main.tex directly

\documentclass[11pt]{report}
\usepackage[margin=0.85in]{geometry}
\usepackage{graphicx}

\usepackage{hyperref}

\newcommand{\mediumsize}{\fontsize{15pt}{16pt}\selectfont}

\begin{document}
	\begin{titlepage}
		
		\begin{minipage}[t]{0.90\textwidth}
			\hfill
			\raggedleft
			\textbf{Edwin Varghese} \\
			London, United Kingdom \\
			\today\\
			\href{mailto:edwinshajivarghese@gmail.com}{edwinshajivarghese@gmail.com} \\
			\url{www.linkedin.com/in/edwinvarghese1}
		\end{minipage}
		
		\raggedright \textbf{Professor Pedro Teixeira-Dias} \\ Centre for Particle Physics and Astronomy \\ Royal Holloway, University of London\\ Egham, UK\\ \href{mailto:Pedro.Teixeira-Dias@rhul.ac.uk}{Pedro.Teixeira-Dias@rhul.ac.uk}\\
		
		\vspace{0.9em}
		
		\raggedright Dear Professor Pedro Teixeira-Dias,\\
		\vspace{0.9em}
		
I am a fourth-year MSci Physics student at Royal Holloway, University of London, and, having benefited from your guidance as my academic tutor, I am writing to apply for your PhD project on Higgs boson production measurements with the ATLAS experiment, with a particular interest in the Higgs boson self-coupling. I am eager to work with the new 13.6 TeV Run 3 dataset and contribute to precision Standard Model measurements, which I see as the next logical step for my development as an experimental particle physicist.\\
\vspace{0.9em}


My third-year project on exploring the phenomenological structure of Higgs-boson decays convinced me that experimental particle physics is the field I want to work in. Driven by curiosity, I learnt that the introduction of the Higgs field and the Higgs mechanism explains mass generation and makes precise predictions that ATLAS can measure, but does not explain why the Yukawa couplings have the values they do, nor does it account for the hierarchy of fermion masses, or the absence of additional scalar states. With this theoretical motivation in place, I could not wait to explore how such ideas are pursued in practice, and it was through research placements that my interests were justified.\\
\vspace{0.9em}


During a six-week research placement, I trained neural-network-based classifiers to evaluate a novel MCMC convergence metric, R*, on simulated datasets from the T2K and NOvA experiments. I benchmarked it against the industry standard Gelman-Rubin (R-hat) statistic, showing that R* detects smaller deviations between chains and provides global convergence information where R-hat is insensitive. The project developed my skills in applying and validating machine learning methods for physics analyses and I am eager to apply these skills as part of a team making precise Higgs self-coupling measurements. At UKAEA, I integrated simulation codes for tokamak plasmas by automating RABBIT input file generation and post-processing plots in Python and benchmarking RABBIT against the higher-fidelity NUBEAM. I learnt to organise large simulation outputs and to write well-documented Python scripts that produce consistent results. I then undertook projects directly relevant to collider phenomenology at Royal Holloway.\\
\vspace{0.9em}


I completed a research review of Higgs searches at the LHC and discussed relevant production, decay modes and discovery channels, with emphasis on VBF. My on-going Master’s project, in which I am deriving the cross section for the hadronic production of the W-boson, gave me new insight into the relationship between theory and phenomenology, where I came to appreciate how the formal structure of quantum field theory leads to predictions that can be tested in experiments.\\
\vspace{0.9em}


Beyond technical skills, I have leadership and communication experience as Academic Officer of the RHUL Physics Society, where I ran LaTeX and Python workshops, organised a weekly PhD seminar series, and coordinated student-led conferences. I enjoy teaching, presenting, and collaborating --- qualities that I believe fit well with the collaborative ethos of ATLAS.\\
\vspace{0.9em}


I am motivated to pursue my PhD at Royal Holloway because of the strong reputation of the RHUL ATLAS group. The group’s leadership in Higgs and top-quark physics, and its recognised expertise in statistical methods create an environment in which I can develop as a well-rounded experimental physicist. Having already studied at Royal Holloway, I have experienced its strong collaborative culture, close links to CERN, and, equally as important, the faculty and technical staff who have personally committed to helping me grow as a student. The group’s involvement in Run 3 analyses and in preparation for the HL-LHC upgrades makes RHUL an ideal place for me to contribute meaningfully to precision Higgs measurements.\\
\vspace{0.9em}


My work across machine learning, simulation work and collider projects has justified my decision to build a career in experimental particle physics. I am eager to consolidate my skills as an analyst in a setting where I can work with real ATLAS data, and it will be a pleasure to exchange ideas with and learn from researchers who are leading precision Higgs measurements. I am excited about the prospect of contributing to a search where new physics may genuinely emerge and I look forward to growing into an independent researcher within the ATLAS community.\\
\vspace{2em}
		
		Thank you for considering my application.\\
		\vspace{2em}
		
		\raggedright Yours sincerely,\\
		\vspace{0.9em}
		\textbf{Edwin Varghese}
		
	\end{titlepage}
\end{document}